\chapter{100}

\section{Mailand (IT), im 
September}

Valentina durchschritt den Flur der \emph{Questur} bedeutend
selbstsicherer als beim ersten Mal. Links und rechts gingen dicht
aneinandergereiht etliche Türen ab, die kleine Büroräume erahnen
liessen. In ihrer Erinnerung war der Korridor länger. Sie ging aufrecht, die
Polizistin Ferrari neben ihr her. Zwischendurch schielte Valentina
unauffällig zu ihr hinüber. Es fühlte sich gut an, beschützt zu sein,
und sie musste zugeben, sich ziemlich wichtig vorzukommen. Schon der
Besuch der beiden Polizisten am Vortag hatte sie ausserordentlich
beeindruckt.

Das Büro des Commissario befand sich einen Stock höher. Die
Räumlichkeiten hier oben waren durchweg grosszügiger, der Flur breiter,
die Dielen höher, die Böden edler.
Die Polizistin klopfte kurz und temperamentvoll an Mantovanis Bürotür.

«Pronto,» klang es tief und laut von drinnen.

Mit der Stimme des Commissario kam die Angst als Druckwelle über
Valentina zurück. Was sich gestern wie eine inszenierte Filmszene
angefühlt hatte, würde gleich wieder erschreckend real sein. Was der
Commissario wohl vorhatte?
Beim Eintreten erschrak Valentina. Entgegen ihrer Vorstellung sah sie
vier Personen an einem länglichen Tisch sitzen. Mantovani an der
Oberseite, längs Rossi und Sönmez und den beiden gegenüber Riva.
Valentina hatte einzig den Commissario erwartet. Wenigstens waren ihr
die Gesichter bereits bekannt. Im Raum hing ein dominanter Kaffeeduft.
Mantovani stand sofort auf, ging auf sie zu und reichte ihr die Hand. Sie fühlte sich einladend an. 

«Buongiorno, Signora Ricci, das hat ja prima geklappt,» er wandte sich
an die junge Polizistin, «danke, Claudia, das wäre vorläufig alles. Ich
melde mich, wenn du sie wieder nach Hause bringen kannst.»

«Bitte.»

Die junge Polizistin verliess den Raum.
Mantovani wies Valentina den Platz zwischen sich und Riva zu. Er setzte sich und lächelte sie freundlich an. 
Dann setzte er sich die Brille auf.

«Wie geht es Ihnen heute, Signora?»

«Ganz okay, Commissario, seit sie,» damit deutete sie auf Rossi und
Sönmez, «seit sie bei mir waren, habe ich weniger Angst.»

Valentina legte ihre Hände in den Schoss. Ihr schien aufzufallen, wie
genau Mantovani jede ihrer Bewegungen verfolgte. Sie sah verlegen aus.

«Das ist gut. Auch, dass Sie sich heute frei nehmen konnten,» er blickte
hoch bevor er sich wieder seinen Notizen zuwandte, «wir wollen zwei
Dinge besprechen. Zum einen, mit welchen Massnahmen wir für Ihre
Sicherheit sorgen werden, zum andern, wie Sie uns helfen können, den
Absturz von Giorgia Sala zu verstehen. Wir gehen davon aus, dass es kein
Selbstunfall war. Aus ermittlungstechnischen Gründen kommunizieren wir
aktuell nicht gegen aussen. Darum kann ich Ihnen nichts dazu sagen. Von Signora Salas Eltern haben wir 
die Erlaubnis eingeholt, Ihnen über den medizinischen Zustand der Verletzten Auskunft geben zu dürfen.
Signora Sala hat nach wie vor keine Erinnerung an den Absturz. Nach der
Einschätzung des behandelnden Arztes ist die Wahrscheinlichkeit, dass
ihre Erinnerung zurückkommt möglich, aber keinesfalls sicher, leider.»

Valentinas Augen vergrösserten sich. Sie sah Giorgia vor sich. Wie elend 
sie ausgesehen hatte. In diesem weiss bezogenen Krankenbett. Wie in ihren Körper gesteckte
Plastikschläuche herumgehangen hatten. Doch am Tragischsten
hatte sie Giorgias Unerreichbarkeit empfunden. Wie weit weg sie war. Das
Bild würde sie wohl ewig verfolgen.

«Das, hm, das ist traurig.»

Emma Riva drehte sich auf ihrem Stuhl ganz in Valentinas Richtung. Ihr
Blick verriet Mitgefühl. Mantovani seufzte.

«Ja, da haben Sie recht,» er faltete bedächtig die Hände, bevor er
weiterfuhr, «nun, Signora, wie Ihnen meine Kollegen mitgeteilt haben,
sind Sie aktuell nur wenig gefährdet. Weil wir aber nicht sicher
sind, von wem die beiden Männer geschickt worden sind, sind wir vorsichtig. 
Natürlich werden wir vorläufig für Ihren Schutz sorgen, Signora.»

Valentina schrak hoch. Emma Riva legte ihr die Hand an
den Oberarm. Und Mantovani breitete beschwichtigend seine Arme aus.

«Signora, kein Grund zur Sorge. Sie bekommen ab sofort \emph{«verdeckten
Polizeischutz».} Tag und Nacht. Sie können uns jederzeit anrufen, wenn
Sie sich unwohl fühlen. Eine unserer Schutzpersonen ist immer in
Ihrer Nähe. Und noch einmal, Signora, die Massnahme ist aktuell eine
vorsorgliche. Sie gehen wie gewohnt an die Uni.»

«Ja, okay, aber, aber was, wenn Louis, also Louis Vincent, die Forderungen 
bis zum 15. September nicht erfüllt hat. Was ist dann mit mir, ich, ich habe Angst.»

Valentina war kurz davor, in Tränen auszubrechen.
Die beiden Polizisten und Riva blickten gleichzeitig zu Mantovani. Als
könnten sie es nicht erwarten, dass er die junge Frau beruhigte.

«Wir verstehen Ihre Angst, Signora. Wir werden alles daran setzen, Sie vor Gefahren zu schützen. Darum geht es heute in erster Linie.»

Valentina rutschte auf dem Stuhl nach vorne, nickte mehrmals und
fixierte Mantovani.
Er zog ein Papier aus einer Ablage und hob es vom Tisch. Dann überflog
er es kurz.

„Wir sind auf ihre Hilfe angewiesen, Signora Ricci. Sie wissen, wir
wollen mit Louis Vincent in Kontakt kommen, um eine Eskalation zu
verhindern. Es wäre ein Leichtes für uns, ihn aufzuspüren, jetzt, wo wir
seine Handynummer haben, aber, weil nichts Konkretes gegen ihn vorliegt,
sind wir nicht befugt, ihn polizeilich zu suchen. Selbst ein
Rechtshilfegesuch aus dem Ausland zur Auslieferung von Vincent ist nicht
gegeben. Es liegt aktuell keine Anklage vor. Doch der Mann könnte uns
ein wichtiger Zeuge werden. Nachdem weder seine Familie noch sein
Arbeitgeber seinen Aufenthaltsort kennen, sind sie unsere 
Verbindung zu ihm. Offenbar sind Sie die einzige Person, der er sich
anvertraut hat?»

«Ich, ich vermute schon.»
